\documentclass[11pt]{article}
\setlength{\topmargin}{-0.8in} % usually -0.25in
\addtolength{\textheight}{1.4in} % usually 1.25in
\addtolength{\oddsidemargin}{-0.8in}
\addtolength{\evensidemargin}{-0.8in}
\addtolength{\textwidth}{1.6in} %\setlength{\parindent}{0pt}

\usepackage[T1]{fontenc} % Recommended for better encoding
\renewcommand{\familydefault}{\sfdefault}
\usepackage{helvet}

% macros
\usepackage{amsmath,
amssymb,
mathrsfs,
amsthm,
fancyhdr,
tikz,
multicol,
enumitem
}

\newcommand{\be}{\begin{enumerate}}
 \newcommand{\ee}{\end{enumerate}}


\newcommand{\C}{{\mathbb{C}}}
\newcommand{\R}{{\mathbb{R}}}
\newcommand{\eps}{\epsilon}
\newcommand{\Z}{{\mathbb{Z}}}
\newcommand{\Q}{{\mathbb{Q}}}
\newcommand{\N}{{\mathbb{N}}}



\lhead{{Math 265 Proofs}}
\chead{\large Midterm III} 
\rhead{Spring 2026}
\cfoot{}
\pagestyle{fancy}
\begin{document}
\thispagestyle{fancy}


\medskip
\large
\vspace{.1in}
\begin{tabular}{l@{\hspace{.4in}}l}
Your Name & \\
\framebox(200,30){} &  \\
\end{tabular}

%\bigskip

\vfill
{
\renewcommand{\baselinestretch}{1.8}
\setlength{\tabcolsep}{.2in}
\normalsize
\begin{center}
\begin{tabular}{|c|c|c|}
\hline
Problem&Total Points&\parbox{.8in}{\hfil Score\hfil}\\
\hline
1&20&\\
\hline
2&15&\\
\hline
3&15&\\
\hline
4&32&\\
\hline
5&18&\\
\hline
\hline
%\hline
Total&100&\\
\hline

\end{tabular}

\end{center}
}
\vfill
\begin{itemize}
\item 
You have 1 hour.

\item If you have a cell phone with you, it should be turned off and put away. (Not in your pocket)

\item You may not use a calculator, book, notes or aids other than a single 3 by 5 notecard.

\item In order to lower the writing load, you may use the symbols $\in$, $\exists$, and $\forall$ \emph{if} they are used precisely and literally. You will not be graded on aesthetics, so repeated use of "Thus, ..." will not be penalized. 
\item You may \emph{not} use the symbol $\Rightarrow$ except parenthetically to indicate the logical structure of an \emph{if-and-only-if} argument. 
\item You must write in complete sentences with proper punctuation. Every sentence must begin with a capitalized word in English, not a symbol. 
\item If you are going to use a method \emph{other than a direct proof}, you must state this at the beginning of your proof unless the directions require you to use particular approach.
\item When using proof by induction, you \emph{must} indicate when the inductive hypothesis is used.
\end{itemize}
\newpage
\begin{enumerate}
%Set Problem
\item (20 pts) 
	\begin{enumerate}
	\item (5 pts) Describe the standard method to prove $X \subseteq Y$ for sets $X$ and $Y$.
	\vspace{1in}
	\item (15 pts) Use the standard method for showing set containment to prove that for all sets $A,$ $B$, and $C$, if $$A\cap C \subseteq B \cap C \text{   and   } A\cup C \subseteq B \cup C,$$
	then $A \subseteq B.$
	\end{enumerate}
\newpage
\item (15 pts) \textbf{Use mathematical induction} to prove $$3^1+3^2+3^3 + \cdots + 3^n = \frac{3^{n+1}-3}{2}$$ for all integers $n \in \N.$
\newpage
\item (15 pts) \textbf{Use mathematical induction} to prove $2^n < (n+1)!$ for all integers $n \geq 2.$
\newpage
\item (8 pts. each for 32 pts. total) Disprove the statements below.
	\begin{enumerate}
	\item If $X \subseteq \N$ and $|X|$ is infinite, then $|\N-X|$ is finite.
	\vfill
	\item For all sets $A$, $B$, and $C,$ if $A \subseteq B,$ then $A \cap \overline{(B \cap C)}=\emptyset.$
	\vfill
	\item For every $n \in \N$, there exists some $m \in \N$ such that $mn \equiv 1 \pmod {10}.$
	\vfill
	\item The relation $R$ on $\Z$ defined by $xRy$ if $|x-y| <5$ is transitive.
	\vfill
	\end{enumerate}
	\newpage
\item (18 pts.) Let $A=\{1,2,3\}.$ Let $R$ be a relation on $\mathcal{P}(A),$ the power set of $A,$ defined as $X\: R\: Y$ if $X \cap \{1,2\} = Y \cap \{1,2\}.$
	\begin{enumerate}
	\item (2 pts) Show that for $X=\{1\}$ and $Y=\{1,3\},$ by the definition, $X\: R\: Y.$
	\vspace{1in}
	\item (10 pts) Prove that $R$ is an equivalence relation.
	\vfill
	\vfill
	\item (6 pts) Describe the partition of $\mathcal{P}(A)$ resulting from the equivalence classes of $R.$ Use proper notation and clearly indicate the members of each partition.
	\vfill
	\end{enumerate}
\end{enumerate}
\end{document}

Let $R$ be a relation on $\N \times \N$ defined as $$(a,b) R (c,d) \text{  if  }a+d=b+c.$$
	\begin{enumerate}
	\item Show that $(1,2)R(3,4).$
	\vfill
	\item Find three additional ordered pairs $(c,d)$ such that $(1,2)R(c,d).$ (It is sufficient to state them. You don't need to prove your are correct.)
	\vfill
	\item Show that $R$ is reflexive.
	\vfill
	\item Show that $R$ is symmetric.
	\vfill
	\item The relation $R$ is an equivalence relation. Determine $[(1,1)]$,  the equivalence class of $(1,1).$
	\vfill
	\end{enumerate}